\section{Conceptual framework}
\label{sec:framework}

Each asset is a county--crop--season profit observation.
Farmers or cooperatives choose crop allocation weights subject to land availability,
crop share bounds, and feasible crop--region combinations.

\subsection{Benchmark strategies}

We compare six benchmarks aligned with the \copper{} charter:
\begin{itemize}
  \item \textbf{B1} observed historical allocation (national area-share proxy where needed);
  \item \textbf{B2} profit-maximizing independent farmer;
  \item \textbf{B3} diversified independent farmer without cross-farmer coordination;
  \item \textbf{B4} local cooperative (primary policy-relevant model);
  \item \textbf{B5} regional cooperative;
  \item \textbf{B6} national cooperative benchmark (upper bound on spatial diversification).
\end{itemize}

\subsection{Objectives and welfare}

Portfolio objectives include expected profit, mean--variance trade-offs, and CVaR-based
downside-risk measures.
Cooperative strategies are evaluated with member-level sharing rules---individual retention,
equal sharing, land-weighted sharing, risk-adjusted sharing, and compensation schemes---and
inequality metrics including Gini coefficients and income dispersion.

Aggregate cooperative gains that harm identifiable members must be flagged; cooperative
performance is not judged on mean profit alone.
